% Chapter 1, Topic from _Linear Algebra_ Jim Hefferon
%  http://joshua.smcvt.edu/linearalgebra
%  2001-Jun-09
\topic{计算的精度}
\index{accuracy! of Gauss's Method|(}
\index{Gauss's Method! accuracy|(}
高斯消元法很适合用计算机来实现。
下面的代码说明了这一点。
它作用于名为~\lstinline[style=inline]!a! 的 $\nbyn{n}$ 矩阵，
先用第一行做倍加，再用第二行，依此类推。
\begin{lstlisting}
for(row=1; row<=n-1; row++){
  for(row_below=row+1; row_below<=n; row_below++){
    multiplier=a[row_below,row]/a[row,row];
    for(col=row; col<=n; col++){
      a[row_below,col]-=multiplier*a[row,col];
    }
  }
}
\end{lstlisting}
这是用 C~语言写的。\index{C language}
这个
\lstinline[style=inline]!for(row=1; row<=n-1; row++){ .. }!
循环把 \lstinline[style=inline]!row! 初始化为~$1$，然后只要
\lstinline[style=inline]!row! 小于等于 $n-1$ 就反复迭代，每次通过
\lstinline[style=inline]!++! 操作把 \lstinline[style=inline]!row! 加一。
另一个不明显的语言结构是，最内层循环中的 \lstinline[style=inline]!-=!
具有如下效果：
\lstinline[style=inline]!a[row_below,col]=-1*multiplier*a[row,col]+a[row_below,col]!。

这段代码虽是高斯消元法机械化的初步尝试，
但它是朴素的。
首先，它假定
\lstinline[style=inline]!row,row!
位置上的元素非零。
因此需要对它做的一种扩充是覆盖这样的情形：在该位置发现零元素时，
要进行对换两行，或者得出矩阵是奇异的结论。

我们可以添加一些 \lstinline[style=inline]!if! 语句来覆盖那些情形，
但我们转而考虑这段代码朴素的另一个方面。
它容易陷入因计算机依赖浮点算术而产生的陷阱。

例如，前面我们已经看到必须把奇异方程组作为单独情形来处理。
但近乎奇异的方程组同样需要小心。
考虑下面这个方程组（多出的数字在第九个有效位上）。
\begin{equation*}
   \begin{linsys}{2}
                   x &+ &2y &= &3\hfill\hbox{}  \\
     1.000\,000\,01x &+ &2y &= &3.000\,000\,01
   \end{linsys}
   \tag{$*$}
\end{equation*}
用肉眼我们很容易看出解为 $x=1$，$y=1$。
计算机则要麻烦得多。
若它用八个有效位表示实数，这称为
\definend{单精度}\index{single precision}，
则它在内部会把第二个方程表示为 $1.000\,000\,0x+2y=3.000\,000\,0$，
丢掉了第九位上的数字。
这台计算机不会报告正确的解，
而会认为两个方程是相同的，
并报告该方程组是奇异的。

为了对计算机怎么会得出相差如此之远的结果有一点直观感受，
考虑该方程组的图像。
\begin{center}
  \includegraphics{gr/mp/ch1.33}
\end{center}
我们无法分辨这两条直线；
该方程组近乎奇异，意思是
这两条直线几乎是同一条直线。
这使得方程组~($*$) 具有如下性质：方程中一个微小的变化
可能引起解的巨大变化。
例如，把 $3.000\,000\,01$ 改为 $3.000\,000\,03$
会使交点从 $(1,1)$ 变到 $(3,0)$。
解会随第九位数字的不同而剧烈变化，
这解释了八位计算机为何会有麻烦。
对输入值的不精确或不确定非常敏感的问题称为
\definend{病态的}。\index{ill-conditioned problem}

上面的例子给出了方程组在计算机上难以求解的一种方式。
它的优点是，两条几乎相同的直线的图像
让人对数值困难发生的一种方式获得难忘的直观认识。
遗憾的是，当我们想求解某个大型方程组时，
这种直观认识没有用处。
对于任意的大型方程组，我们通常并不了解其几何。

除了某些线性曲面之间的夹角很小之外，
计算机的结果不可靠还有其他原因。
例如，考虑这个方程组（来自 \cite{Hamming}）。
\begin{equation*}
  \begin{linsys}{2}
     0.001x  &+  &y  &=  &1  \\
          x  &-  &y  &=  &0
  \end{linsys}
\tag*{($**$)}\end{equation*}
第二个方程
给出 $x=y$，于是 $x=y=1/1.001$，
因而两个变量的值都恰好略小于 $1$。
使用两位数字的计算机在内部如下表示该方程组
（为清楚起见，我们用两位浮点
算术来做这个例子，但构造一个用八位或更多位的类似例子并不难）。
\begin{equation*}
  \begin{linsys}{2}
    (1.0\times 10^{-3})\cdot x  &+  &(1.0\times 10^{0})\cdot y  &=  &1.0\times 10^{0}  \\
    (1.0\times 10^{0})\cdot x   &-  &(1.0\times 10^{0})\cdot y  &=  &0.0\times 10^{0}
  \end{linsys}
\end{equation*}
行化简步骤 $-1000\rho_1+\rho_2$ 产生第二个方程 $-1001y=-1000$，
这台计算机把它舍入到两位，得 $(-1.0\times 10^{3})y=-1.0\times 10^{3}$。
计算机由第二个方程确定 $y=1$，
并据此从第一个方程得出 $x=0$。
$y$~值接近，但 $x$~值很糟糕\Dash
实际答案与计算机答案之比是
无穷大。\index{accuracy!rounding error}
简言之，当求解方程组的代码导致使用较小的首主元时，
计算机对浮点算术的依赖是输出不可靠的另一个原因。

有经验的程序员可能会以使用
\definend{双精度}来应对，\index{double precision} %
它保留十六位有效数字，
或者使用更大的位数。
这确实能解决许多问题。
然而，双精度有更大的内存需求，
而且显然我们可以照样调整上面的例子，
使同样的麻烦出现在第十七位上，所以双精度并非万灵药。
我们需要一个使数值麻烦最小化的策略，
还需要一些关于所报告的解可以信任到什么程度的指导。

对上面朴素代码的一个基本改进是：不要简单地取
\lstinline[style=inline]!row,row! 位置上的元素
来作为倍加所用的因子，
而是查看 \lstinline[style=inline]!row,row! 元素下方
\lstinline[style=inline]!row! 列中的全部元素，
并选取一个因为不太小而可能给出可靠结果的元素。
这就是\definend{部分主元法}。%
\index{partial pivoting}\index{pivoting!full} 

例如，要求解上面这个棘手的方程组~($**$)，
我们先查看两个方程以寻找最适合使用的元素，
并取第二个方程中的~$1$，因为它更可能给出好结果。
倍加步骤 $-.001\rho_2+\rho_1$ 给出第一个方程
$1.001y=1$，计算机会把它表示为
$(1.0\times 10^{0})y=1.0\times 10^{0}$，
从而得出 $y=1$，再经回代得出 $x=1$，
两者都接近正确。
我们可以把上面的代码改写来做这件事。
\begin{lstlisting}
for(row=1; row<=n-1; row++){
/* find the largest entry in this column (in row max) */
  max=row;
  for(row_below=row+1; row_below<=n; row_below++){
    if (abs(a[row_below,row]) > abs(a[max,row]));
      max = row_below;
  }
/* swap rows to move that best entry up */
  for(col=row; col<=n; col++){
    temp=a[row,col];
    a[row,col]=a[max,col];
    a[max,col]=temp;
  }
/* proceed as before */
  for(row_below=row+1; row_below<=n; row_below++){
    multiplier=a[row_below,row]/a[row,row];
     for(col=row; col<=n; col++){
       a[row_below,col]-=multiplier*a[row,col];
    }
  }
}
\end{lstlisting}

对实现高斯消元法的最佳方式的完整分析超出了本书的范围
（见 \cite{Wilkinson65}），
但大多数专家推荐的方法
先在候选元素中找出最好的元素，
然后把它缩放成一个不太可能带来麻烦的数。
这就是
\definend{比例部分主元法}。\index{scaled partial pivoting}\index{pivoting!partial!scaled}

除了返回可能可靠的结果之外，
多数完善的代码还会返回一个
\definend{条件数}%
\index{conditioning number}\index{matrix!conditioning number}
它描述输入数字中的不确定
可能被放大而成为返回结果中不精确的倍数
（见 \cite{Rice}）。

教训是：
高斯消元法在理论上总是有效，
计算机代码也正确地实现了该方法，
并不意味着答案就是可靠的。
实践中，总要使用这样的软件包：
其中的专家已经努力防范各种可能出问题的地方。

\begin{exercises}
  \item 
    用两位小数，把 $253$ 与 $2/3$ 相加。
    \begin{answer}
      科学记数法便于表达两位的限制：
      $.25\times 10^{2}+.67\times 10^{0}=.25\times 10^{2}$。
      注意加上 $2/3$ 对总和没有影响。
    \end{answer}
  \item 
    这个直线相交问题与上面讨论的例子形成对比。
    \begin{center}
      \vcenteredhbox{\includegraphics{gr/mp/ch1.34}}
      \qquad
      $\displaystyle \begin{linsys}{2}
            x &+ &2y &= &3  \\
            3x &- &2y &= &1
      \end{linsys}$
    \end{center}
    把底下方程中的常数改为 $1.008$ 并求解，
    以此说明：该方程组中数字的某个小变化
    只会产生解的小变化。
    并将它与未改变的方程组的解相比较。
    \begin{answer}
      消元
      \begin{equation*}
        \grstep{-3\rho_1+\rho_2}
        \begin{linsys}{2}
          x  &+  &2y   &=  &3  \\
             &   &-8y  &=  &-7.992
        \end{linsys}
      \end{equation*}
      给出 \( (x,y)=(1.002,0.999) \)。
      所以对该方程组而言，常数的微小变化只产生解的微小变化。
    \end{answer}
  \item 
    考虑这个方程组（\cite{Rice}）。
    \begin{equation*}
      \begin{linsys}{2}
        0.000\,3x  &+  &1.556y  &=  &1.559 \\
        0.345\,4x  &-  &2.346y  &=  &1.108
      \end{linsys}
    \end{equation*}
    \begin{exparts*}
      \partsitem 求解它。
      \partsitem 用每一步都舍入到四位数字来求解它。
    \end{exparts*}
    \begin{answer}
      \begin{exparts}
        \partsitem 完全精确的解是 $x=10$ 与 $y=1$。
        \partsitem 四位数字的消元
          \begin{equation*}
            \grstep{-(.3454/.0003)\rho_1+\rho_2}
            \begin{amat}{2}
              .0003  &1.556  &1.569  \\
              0      &-1794   &-1805
            \end{amat}
          \end{equation*}
          给出结论~$y=1.006$，这还不错，
          以及 $x=12.21$。
          当然，这与正确答案相差百分之二十。
      \end{exparts}
    \end{answer}
  \item 
    计算机内部的舍入常常对结果有影响。
    假设你的机器有八位有效数字。
    \begin{exparts}
      \partsitem 证明机器会算出
         $(2/3)+((2/3)-(1/3))$ 不等于 $((2/3)+(2/3))-(1/3)$。
         因此，计算机算术不满足结合律。
      \partsitem 比较计算机所表示的 $(1/3)x+y=0$
         与 $(2/3)x+2y=0$。
         第一个方程的两倍与第二个方程相同吗？
    \end{exparts}
    \begin{answer}
      \begin{exparts}
        \partsitem 对于第一个，首先 $(2/3)-(1/3)$ 为
          $.666\,666\,67-.333\,333\,33=.333\,333\,34$，
          于是
          $(2/3)+((2/3)-(1/3))=.666\,666\,67+.333\,333\,34=1.000\,000\,0$。

          对于另一个，首先
          $((2/3)+(2/3))=.666\,666\,67+.666\,666\,67=1.333\,333\,3$，
          于是
          $((2/3)+(2/3))-(1/3)=1.333\,333\,3-.333\,333\,33=.999\,999\,97$。
        \partsitem 第一个方程是
          $.333\,333\,33\cdot x+1.000\,000\,0\cdot y=0$，
          而第二个方程是
          $.666\,666\,67\cdot x+2.000\,000\,0\cdot y=0$。
      \end{exparts}
    \end{answer}
  \item 
    病态不仅依赖于系数矩阵。
    这个例子 \cite{Hamming} 表明，
    病态也可能由方程组左边与右边之间的相互作用产生。
    设 $\varepsilon$ 是一个小的实数。
    \begin{equation*}
      \begin{linsys}{3}
        3x  &+  &2y           &+  &z            &=  &6\hfill   \\
        2x  &+  &2\varepsilon y  &+  &2\varepsilon z  
                                          &=  &2+4\varepsilon\hfill \\
         x  &+  &2\varepsilon y  &-  &\varepsilon z   
                                          &=  &1+\varepsilon\hfill
      \end{linsys}
    \end{equation*}
    \begin{exparts}
      \partsitem 用手算求解该方程组。
        注意，$\varepsilon$ 之所以能约去，只是因为
        右边与左边一样，整数部分恰好完全消去。
      \partsitem 用手算求解该方程组，舍入到两位
        小数，并取 $\varepsilon=0.001$。
%      \partsitem If you have access to a computer package for solving
%        linear systems, see if you can find an $\varepsilon$ small enough
%        to get your package to give incorrect results.
    \end{exparts}
    \begin{answer}
      \begin{exparts}
        \partsitem 这个计算
          \begin{align*}
            &\grstep[-(1/3)\rho_1+\rho_3]{-(2/3)\rho_1+\rho_2}
            \begin{amat}{3}
              3  &2                   &1                   &6               \\
              0  &-(4/3)+2\varepsilon &-(2/3)+2\varepsilon &-2+4\varepsilon \\
              0  &-(2/3)+2\varepsilon &-(1/3)-\varepsilon  &-1+\varepsilon 
            \end{amat}                                                    \\
            &\grstep{-(1/2)\rho_2+\rho_3}
            \begin{amat}{3}
              3  &2                   &1                   &6               \\
              0  &-(4/3)+2\varepsilon &-(2/3)+2\varepsilon &-2+4\varepsilon \\
              0  &\varepsilon         &-2\varepsilon       &-\varepsilon 
            \end{amat}
          \end{align*}
          给出第三个方程 $y-2z=-1$。
          代入第二个方程给出
          $((-10/3)+6\varepsilon)\cdot z=(-10/3)+6\varepsilon$，
          所以 $z=1$，从而 $y=1$。
          有了这些，第一个方程给出 $x=1$。
        \partsitem
          与上面一样，科学记数法便于表达
          对位数的限制。

          保留两位数字的解
          为 $z=2.1$，$y=2.6$，$x=-.43$。
          \begin{multline*}
            \begin{amat}{3}
              .30\times 10^{1}  &.20\times 10^{1}  &.10\times 10^{1} 
                 &.60\times 10^{1}        \\
              .10\times 10^{1}  &.20\times 10^{-3} &.20\times 10^{-3} 
                 &.20\times 10^{1}        \\
              .30\times 10^{1}  &.20\times 10^{-3}  &-.10\times 10^{-3} 
                 &.10\times 10^{1}        
            \end{amat}                                           \\
            \begin{aligned}
            &\grstep[-(1/3)\rho_1+\rho_3]{-(2/3)\rho_1+\rho_2}
            \begin{amat}{3}
              .30\times 10^{1}  &.20\times 10^{1}  &.10\times 10^{1} 
                 &.60\times 10^{1}        \\
              0                 &-.13\times 10^{1} &-.67\times 10^{0} 
                 &-.20\times 10^{1}        \\
              0                 &-.67\times 10^{0}  &-.33\times 10^{0} 
                 &-.10\times 10^{1}        
            \end{amat}                                          \\
            &\grstep{-(.67/1.3)\rho_2+\rho_3}
            \begin{amat}{3}
              .30\times 10^{1}  &.20\times 10^{1}  &.10\times 10^{1} 
                 &.60\times 10^{1}        \\
              0                 &-.13\times 10^{1} &-.67\times 10^{0} 
                 &-.20\times 10^{1}        \\
              0                 &0                  &.15\times 10^{-2} 
                 &.31\times 10^{-2}        
            \end{amat}
            \end{aligned}
          \end{multline*}
      \end{exparts}
    \end{answer}
\end{exercises}
\index{accuracy! of Gauss's Method|)}
\index{Gauss's Method! accuracy|)}

\endinput



